\section{Statement and context}

Campana's theory of special varieties, developed in \cite{Campana2004}, describes the part of birational geometry with no positive-dimensional general-type orbifold base. The orbifold formulation records the multiplicities of fibres and divisors. The orbifold Abelianity Conjecture predicts virtual abelianity for the fundamental group of a smooth integral special geometric orbifold whose coarse space is bimeromorphic to a compact K\"ahler manifold; see \cite[Conjecture~12.10(1)]{CampanaOrbifolds} and its projective formulation in \cite[Conjecture~11.2]{CampanaSurvey}. Here, as there, virtually abelian means containing an abelian subgroup of finite index.

We give a conditional resolution of the case consisting of a smooth projective fourfold and one smooth divisor of multiplicity two. The argument uses the manifold abelianity theorem in every dimension and constructs a special manifold whose fundamental group surjects onto the orbifold group.

We first fix the differential definition of specialness used throughout. Let $X$ be a connected smooth projective complex variety and $D$ a smooth reduced divisor. Its order-two root orbifold is denoted by
\[
 \pi:\cX=\sqrt[2]{(X,D)}\longrightarrow X.
\]
Near $D=\{z_1=0\}$ it is presented by
\begin{equation}\label{eq:root-chart}
 z_1=w_1^2,\qquad z_i=w_i\ (i>1),
 \qquad w_1\longmapsto-w_1.
\end{equation}
Bundles, forms and their maps on $\cX$ mean equivariant holomorphic objects in these smooth charts, with the compatible identifications on overlaps. Write $\cD$ for the reduced root divisor. Its divisor line has the canonical relation
\begin{equation}\label{eq:root-line}
 T:=\OO_{\cX}(\cD),\qquad T^{\otimes2}=\pi^*\OO_X(D).
\end{equation}
The stabilizer acts by sign on $T$ along $\cD$. This line exists on $\cX$ even when $\OO_X(D)$ has no square root on $X$.

For an ordinary or orbifold line bundle $L$, its Iitaka dimension $\kappa(L)$ is $-\infty$ if every positive power has no nonzero section, and otherwise is the largest dimension of the images of its complete positive-power section maps. On an orbifold these are the global equivariant sections.

\begin{definition}\label{def:special}
A smooth compact complex manifold or a smooth root orbifold $V$ is \emph{special} if there is no integer $1\le p\le\dim V$ and no holomorphic line bundle $L$ admitting a nonzero sheaf map
\[
 L\longrightarrow\Omega_V^p,\qquad \kappa(L)=p.
\]
\end{definition}

This is the finite-integral root form of the Bogomolov-sheaf criterion in \cite[Theorem~8.9]{CampanaOrbifolds}. We retain Definition~\ref{def:special} throughout the proof. Remark~\ref{rem:conventions} gives the local comparison with orbifold symmetric differentials.

\begin{assumption}[Manifold abelianity]\label{ass:manifold}
Every connected smooth compact K\"ahler manifold which is special in the sense of Definition~\ref{def:special} has virtually abelian entire discrete topological fundamental group, in every complex dimension.
\end{assumption}

This is Theorem~1.1 of the companion paper \cite{Supplied}. No orbifold extension of that theorem is assumed. The following theorem, including the appendix proof, uses no other result of that paper.

\begin{theorem}\label{thm:main}
Assume Assumption~\ref{ass:manifold}. Let $X$ be a connected smooth projective complex fourfold and let $D\subset X$ be a nonempty smooth connected reduced divisor. If $\sqrt[2]{(X,D)}$ is special, then
\begin{equation}\label{eq:target-group}
 G=\pi_1(X\setminus D,x)\big/\langle\!\langle\gamma^2\rangle\!\rangle
\end{equation}
contains an abelian subgroup of finite index. Here $\gamma$ is a positively oriented meridian about $D$, transported to $x\in X\setminus D$.
\end{theorem}

The assertion concerns the entire discrete group in \eqref{eq:target-group}, not just its abelianization, finite quotients or linear images. No uniform index is claimed, and torsion is allowed. We impose no residual-finiteness or linearity assumption on $G$, no positivity condition on $K_X+\tfrac12D$, and no finite manifold-cover hypothesis.

The main construction is independent of specialness. It produces a smooth projective eightfold $Y$ over $\cX$ with general fibre a product of four genus-one curves. It also supplies a submersive lift at some point over the general point of \emph{every} base divisor. These points detect all poles in the descent of differential tensors. A direct argument then proves that specialness of $\cX$ forces specialness of $Y$, while even pullback multiplicities give an epimorphism $\pi_1(Y)\twoheadrightarrow G$. These two conclusions permit the application of Assumption~\ref{ass:manifold}.

Proposition~\ref{prop:transfer} isolates this passage in arbitrary base dimension for an order-two smooth boundary. Its hypotheses concern general fibres and divisorial submersions, so it does not require a globally smooth family. In the fourfold construction, four independent genus-one factors eliminate stabilizers over the three-dimensional divisor. Appendix~\ref{app:branched} gives an alternative construction using auxiliary branched covers, affine elliptic actions and a quotient resolution. Both constructions avoid a finite uniformization requirement.

The alternative proof \emph{Conditional orbifold abelianity by inertia and quotient minimality} retains the longer descent method and states its additional inputs from the companion paper. It is not needed for the proof of Theorem~\ref{thm:main} below.

\paragraph{Conventions.}
All varieties are over $\CC$. Projective spaces of vector bundles parametrize lines. A general fibre is taken over a sufficiently small nonempty Zariski open. Fundamental-group maps use compatible base points and paths; subgroup images are understood up to conjugacy. We use the standard complex-projective comparison between algebraic and analytic bundles, sections, and finite coherent algebras \cite{GAGA}.
